% !TEX program = xelatex
% Upload the entire ZIP to Overleaf. Set the compiler to XeLaTeX.
% Template faithful to the typography of chrelli/bioRxiv-word-template.
% The bundled vector bioRxiv logo is inserted by biorxiv-chrelli.cls.
% No separate upload or SVG conversion is required in Overleaf.
\documentclass{biorxiv-chrelli}

\ShowBioRxivLogo % Comment out to hide logo AND remove its spacing

\usepackage{lipsum} % Only for demonstration; remove for your manuscript.

% --- EDIT THE PAPER METADATA HERE ---
\PaperTitle{Challenging experiments reveal a very surprising
\\and interesting biological discovery}
\PaperAuthors{First A. Lastname\affilmarker{1}, Second B. Lastname\affilmarker{2,3}
  and First B. Lastname\affilmarker{1,4,*}}
\PaperAffiliations{%
  \affilmarker{1}Institute of Biology, University of Something, 10115 Berlin, Germany\\
  \affilmarker{2}Center for Biological Sciences, Another University, 10115 Berlin, Germany\\
  \affilmarker{3}Current address: Department of Animal Biology, University Somewhere Else, NY, USA\\
  \affilmarker{4}Biology Cluster of Whatever, Some Research University, 2200 Copenhagen N, Denmark\\
  \affilmarker{*}Correspondence should be addressed to F.B.L. (first.lastname@example.edu)
}
\PaperAbstract{%
  Biological research often benefits from clear, compact manuscript design.
  This example reproduces the principal styles of the Chrelli bioRxiv Word
  template: a Helvetica-like title and abstract, STIX Two body and mathematics,
  red headings, narrow two-column text, and carefully sized figure legends.
  Replace this abstract with your own scientific summary. A well-written
  abstract states the question, gives the main experimental approach, reports
  the principal observations and concludes with their wider significance.
  This is demonstration text rather than a statement of scientific findings.
  The class is an independent LaTeX adaptation intended for an editable
  Overleaf manuscript, not a publishing service or an official bioRxiv style.
  \lipsum[1] \lipsum[2][1-5]
}
\RunningTitle{Running title}
\FooterCitation{Lastname et al. 2026 (preprint)}

\begin{document}
\maketitle

\section{Introduction}
\lipsum[1-3]
\lipsum[4][1-6]

% Full-width or 1.5-column figures use figure*.
% In LaTeX these are top-of-page floats; exact Word-style image anchoring
% is intentionally replaced by reliable LaTeX float placement.
\begin{figure*}[t]
  \centering
  \begin{minipage}{133.125mm}
    \centering
    \colorbox{biorxivblue}{\parbox[c][26mm][c]{\dimexpr\linewidth-2\fboxsep\relax}{\strut}}
    \caption{\legendtitle{Example of a one-and-one-half-column figure.}
      This 133.125\,mm-wide environment uses the same figure width recommended
      by the source Word template. It may sit at the top of a following page.}
    \label{fig:onehalf}
  \end{minipage}
\end{figure*}


\section{Results}
\ResultsLead{A short red lead-in introduces the experimental result}
\lipsum[5]
\lipsum[6][1-6]

\begin{figure}[t]
  \centering
  % Original Word example uses a 89 mm one-column figure.
  \colorbox{biorxivblue}{\parbox[c][43mm][c]{\dimexpr\columnwidth-2\fboxsep\relax}{\strut}}
  \caption{\legendtitle{Example of a one-column figure.}
  The blue rectangle stands in for an experimental image or graph.
  The figure occupies exactly one 89\,mm text column. Replace this
  placeholder with \texttt{\textbackslash includegraphics[width=\textbackslash columnwidth]\{myfigure\}}.}
  \label{fig:one}
\end{figure}

\subsection{Another example of a subheading}
\lipsum[7-8]

\section{Discussion}
\lipsum[9]\lipsum[10]
\ResultsLead{What a massive breakthrough for the field}
\lipsum[11]

% A full-width figure spans the full 183.9 mm text area (both columns).
% Like the 1.5-column example, figure* is a top-of-page float. In this
% three-page mockup, it demonstrates the third image size on page 3.
\begin{figure*}[t]
  \centering
  \colorbox{biorxivblue}{\parbox[c][39mm][c]{\dimexpr\textwidth-2\fboxsep\relax}{\strut}}
  \caption{\legendtitle{Example of a full-width, two-column figure.}
    This 183.9\,mm-wide placeholder spans the entire text area and illustrates
    the largest figure size in the template. Replace it with
    \texttt{\textbackslash includegraphics[width=\textbackslash textwidth]\{myfigure\}}.}
  \label{fig:fullwidth}
\end{figure*}

% Additional sample paragraphs help demonstrate the three-page layout.
\lipsum[15]

\section{Materials and Methods}
{\MethodsText
\MethodsLead{First methods subheading (no space before)}
\lipsum[12]
\MethodsLead{Very tricky method (4 pt space before)}
\lipsum[13-14]
\lipsum[16-23]
}

\section{References and citations}
A citation in the same compact authorial text style \citep{example2026}
can be managed with your BibTeX database. Equations use STIX Two Math:
\begin{equation}
  P(Y\!=\!1\mid X)=\frac{e^{\beta_0+\beta_1X}}{1+e^{\beta_0+\beta_1X}}.
\end{equation}

% In your real paper, remove the example section above and leave just
% \bibliography{references} so the bibliography begins with "References".
\bibliographystyle{unsrtnat}
\bibliography{references}

\section{Acknowledgements}
We thank our colleagues and collaborators for constructive discussion.
Replace this text with your acknowledgements and funding information.

\section{Author contributions}
Conceptualization: F.A.L., S.B.L.; Investigation: F.A.L.; Methodology: S.B.L.;
Writing--original draft: F.A.L.; Writing--review and editing: all authors.

\section{Competing interest statement}
The authors declare no competing interests.

\end{document}
